\chapter{Excessive Sweating}

\section*{Definition}
Excessive sweating (hyperhidrosis) is sweating beyond thermoregulatory needs that causes discomfort or interferes with daily, social, or occupational activities.

\section*{When to See a Doctor}
See your doctor if:
\begin{itemize}
  \item Sweating is new, unexplained, or accompanied by weight loss, fever, or palpitations
\end{itemize}

\section*{How It Develops}
Primary focal hyperhidrosis usually affects the palms, soles, underarms, or face without an identifiable systemic cause. Secondary generalized sweating can result from medicines, endocrine or metabolic disease, infection, menopause, malignancy, or neurological disease.

\section*{Causes}
\begin{itemize}
  \item Primary hyperhidrosis focal excessive sweating
  \item Medications antidepressants opioids
  \item Hyperthyroidism diabetes mellitus
  \item Malignancy lymphoma infections
  \item Menopausal hot flashes withdrawal
\end{itemize}

\section*{Related Symptoms}
\begin{itemize}
  \item Night sweats chills anxiety weight loss
  \item Fever palpitations orthostatic intolerance
  \item Dehydration electrolyte imbalance social isolation embarrassment
  \item Clothing stains occupational limitation reduced quality of life
  \item Social avoidance psychological impact depression anxiety
\end{itemize}

\section*{Medical Evaluation}
\begin{itemize}
  \item History assesses distribution, timing, triggers, sleep-related sweating, medicines, substances, menopause, and associated symptoms.
  \item Examination and selected tests are guided by whether primary focal or secondary generalized hyperhidrosis is suspected.
  \item Tests may include blood counts, glucose, thyroid studies, or infection evaluation when indicated; extensive testing is not routine for typical focal disease.
\end{itemize}

\section*{Treatment Options}
\begin{itemize}
  \item Aluminum chloride antiperspirant is often tried first for primary focal hyperhidrosis.
  \item Prescription topical or oral medicines, botulinum toxin injections, iontophoresis, or specialist procedures may be considered when needed.
  \item Secondary sweating is treated by addressing the underlying condition or causative medicine with clinical guidance.
\end{itemize}

\section*{Self-Care and Home Management}
\begin{itemize}
  \item Wear breathable clothing and use absorbent pads or socks when helpful.
  \item Apply antiperspirant to completely dry skin as directed; irritation may require clinician advice.
  \item Replace fluids after heavy sweating and protect skin from maceration or irritation.
  \item Keep a record of triggers, medicines, timing, and whether sweating occurs during sleep.
\end{itemize}

\section*{References}
\begin{thebibliography}{99}
\bibitem{excessive_sweating:ref1} Hornberger J, Grimes K, Naumann M, et al. ``Recognition, Diagnosis, and Treatment of Primary Focal Hyperhidrosis.'' J Am Acad Dermatol. 2004;51(2):274--286.
\bibitem{excessive_sweating:ref2} Solish N, Bertucci V, Dansereau A, et al. ``A Comprehensive Approach to the Recognition, Diagnosis, and Severity-Based Treatment of Focal Hyperhidrosis.'' J Am Acad Dermatol. 2007;56(2):303--310.
\bibitem{excessive_sweating:ref3} Haider A, Solish N. ``Focal Hyperhidrosis: Diagnosis and Management.'' CMAJ. 2005;172(1):69--75.
\bibitem{excessive_sweating:ref4} Walling HW. ``Primary Hyperhidrosis: A Review.'' Dermatol Surg. 2011;37(5):537--554.
\bibitem{excessive_sweating:ref5} Nawrocki S, Cha J. ``The Etiology, Diagnosis, and Management of Hyperhidrosis: A Comprehensive Review.'' J Am Acad Dermatol. 2019;81(3):657--667.
\end{thebibliography}
